\section{Подготовка к устной защите}
\label{sec:defense}

\subsection{Как последовательно объяснить свой регулятор}

Для каждого пункта работы полезно восстановить одно связное доказательство.
Начните с модели, измеряемых переменных и целевого равновесия. Затем
объясните, какая идея метода определила замену координат или поверхность.
Выведите управление из нужной динамики и только после этого сформулируйте
результат об устойчивости. В конце назовите область определения закона,
необходимые предположения и смысл численного опыта.

На защите важно уметь переходить в обе стороны: из формулы закона получить
уравнение замкнутой системы и из желаемого уравнения снова вывести закон.
Это сразу выявляет потерянный знак, пропущенную производную или неизвестный
параметр, случайно попавший в реализуемый регулятор.

\subsection{Общие вопросы}

\begin{enumerate}
\item \textbf{Чем устойчивость отличается от притяжения?}
Устойчивость удерживает близкие начальные состояния близко к равновесию.
Притяжение означает сходимость к равновесию. Асимптотическая устойчивость
требует обоих свойств; экспоненциальная дополнительно задаёт скорость
$\norm{x(t)}\le M e^{-\lambda(t-t_0)}\norm{x(t_0)}$.

\item \textbf{Почему $\dot V\le0$ недостаточно для асимптотической устойчивости?}
Вне нуля могут существовать траектории с постоянной $V$. Нужно исследовать
наибольшее инвариантное множество в $\{\dot V=0\}$, например по ЛаСаллю.
При прямой равновесий получается устойчивость множества, но не притяжение
к одной выбранной точке этой прямой.

\item \textbf{Что добавляет радиальная неограниченность?}
Она делает подуровни $V$ ограниченными и позволяет переносить локальное
рассуждение об убывании на любые начальные состояния, если условия на
производную выполняются глобально. Также нужна определённость $V$
относительно правильного равновесия и полнота решений.

\item \textbf{Гурвицева матрица всегда даёт глобальный результат?}
Для линейной системы --- да. Если это только якобиан нелинейной системы
в равновесии, результат локальный. Глобальный вывод после точной
линеаризации требует глобальной замены координат и глобально определённой
обратной связи.

\item \textbf{Сохраняется ли глобальная экспоненциальная оценка при любой замене координат?}
Глобальная асимптотическая устойчивость сохраняется при глобальном
диффеоморфизме, переводящем равновесие в нуль. Для той же глобальной
экспоненциальной оценки по исходной норме нужны дополнительные глобальные
сравнения норм, например двусторонние липшицевы оценки преобразования.

\item \textbf{Чем отличаются finite-time и fixed-time?}
При finite-time время достижения конечно для каждого начального состояния,
но может расти с его величиной. При fixed-time имеется общая верхняя граница
времени для всех начальных состояний. Конечное время достижения поверхности
скольжения не означает конечного времени достижения начала координат.
\end{enumerate}

\subsection{Вопросы по работе 3}

\begin{enumerate}[resume]
\item \textbf{Что такое относительная степень $\rho$?}
Это номер первой производной выхода, содержащей управление с ненулевым
коэффициентом. Предыдущие коэффициенты $L_gL_f^{i-1}h$ должны равняться нулю
в рассматриваемой области, а последний $L_gL_f^{\rho-1}h$ --- не обращаться
в нуль. Одной проверки в произвольной отдельно взятой точке недостаточно
для утверждения об области.

\item \textbf{Почему внутренние координаты выбираются из $L_g\varphi_i=0$?}
Тогда $\dot\eta_i=L_f\varphi_i$ не содержит непосредственного слагаемого
с $u$. Но управление по-прежнему влияет на внутреннюю динамику через
внешние координаты $\xi$; внутренние переменные не обязательно автономны
при ненулевом выходе.

\item \textbf{Почему одного ненулевого якобиана недостаточно для глобальности?}
Теорема об обратной функции даёт локальную обратимость.
Разные исходные точки всё ещё могут переходить в одинаковые новые
координаты. Нужны инъективность и подходящий глобальный обратный переход;
в задаче 4 две ветви $x_3>0$ и $x_3<0$ имеют одно и то же изображение.

\item \textbf{Что такое нуль-динамика?}
Это движение внутренних координат при поддержании выхода и его нужных
производных тождественно нулевыми. В нормальной форме она равна
$\dot\eta=f_0(\eta,0)$. При этом вход обычно равен компенсации
$u=\alpha(x)$, а не обязательно нулю.

\item \textbf{Почему система 1 не минимально-фазовая?}
Её нуль-динамика имеет собственные значения $0,-2$ и прямую равновесий
$\eta_1=\eta_2$. Составляющая вдоль этой прямой не затухает.
По лекционному определению нужна асимптотическая устойчивость нуль-динамики.

\item \textbf{Почему система 2 стабилизируется глобально?}
Замена $\eta=x_3-x_1^2/2$, $\xi_1=x_1$, $\xi_2=-x_1+x_2$
глобально обратима. При $u=x_1x_3-x_1$ вся система в новых координатах
линейна с полюсами $-2,-1,-1$. Устойчивость доказана для всех новых
координат, включая внутреннюю $\eta$.

\item \textbf{Как система 3 может иметь нулевой выход и расходящееся состояние?}
При $\xi_1(0)=\xi_2(0)=0$, $\eta(0)\ne0$ внешняя динамика остаётся нулевой,
а $x_3=\eta=\eta(0)e^t$. Линеаризующий вход компенсирует влияние $x_3$
на выход, но не устраняет неустойчивость самого $x_3$.

\item \textbf{Что особенного в системе 4?}
И коэффициент $L_gL_fh=x_3$, и определитель $DT=-x_3$ вырождаются на
$x_3=0$. На каждой ветви есть обратный корень
$x_3=\pm\sqrt{2(\xi_2-\eta)}$. Нуль-динамика приближается к границе
регулярной области, поэтому нельзя назвать преобразование и дробь
$(x_1+v)/x_3$ регулярными в предельной точке.

\item \textbf{Почему в системе 5 $K=(12,16,7)$, а не $(7,16,12)$?}
Для цепочки третьего порядка с $v=-k_1z_1-k_2z_2-k_3z_3$
характеристический полином равен $s^3+k_3s^2+k_2s+k_1$.
Он должен совпасть с $(s+2)^2(s+3)=s^3+7s^2+16s+12$.
После подстановки $z_3=x_3+x_1^2$ обязательно появляется член $-7x_1^2$.
\end{enumerate}

\subsection{Вопросы по работе 4}

\begin{enumerate}[resume]
\item \textbf{Что является виртуальным управлением?}
Это желаемое значение следующего состояния, например $x_2=\phi(x_1)$.
Его нельзя подать непосредственно на объект: физический вход управляет
производной следующего состояния. Поэтому вводится ошибка
$z_2=x_2-\phi$ и строится следующий шаг.

\item \textbf{Откуда берётся $\dot\phi$?}
Из дифференцирования ошибки: $\dot z=\dot\xi-\dot\phi$.
Это полная производная по исходной системе,
$\dot\phi=\sum_j(\partial\phi/\partial x_j)\dot x_j$.
Пропуск этого члена изменяет замкнутую динамику и разрушает сокращение
смешанных слагаемых в $\dot V_c$.

\item \textbf{Почему производные виртуальных управлений ещё являются обратной связью по состоянию?}
В строгой треугольной структуре $\phi_i$ зависит только от первых $i$
координат, производные которых до последнего шага не содержат настоящего
входа $u$. Поэтому $\dot\phi_i$ вычисляется через измеренное состояние.
Это нужно проверить, а не использовать как безусловное свойство любой системы.

\item \textbf{Что удерживает третью систему в ненулевой точке?}
Для неё $x_{ss}=(0,1,0,0)^T$, \mbox{$u_{ss}=0$}. В нуле $\dot x_1=1$,
поэтому нуль не равновесие. В целевой точке $\cos0-x_2=1-1=0$,
а остальные производные также нулевые. Ошибки бэкстеппинга обращаются
в нуль именно в этой точке.

\item \textbf{Почему в производной $\phi_3$ нужен дополнительный член от знаменателя?}
При $\phi_3=N/g_2$, $g_2=2-\sin x_3$,
\[
 \frac{\partial\phi_3}{\partial x_3}
 =\frac{\partial N/\partial x_3}{g_2}
  +\frac{N\cos x_3}{g_2^2}.
\]
Плюс появляется потому, что $\partial g_2/\partial x_3=-\cos x_3$.
Граница $g_2\ge1$ гарантирует отсутствие сингулярности всех шагов.
\end{enumerate}

\subsection{Вопросы по работе 5}

\begin{enumerate}[resume]
\item \textbf{Что выбирается первым --- усиление или поверхность?}
Сначала поверхность, на которой движение должно быть устойчивым, затем
усиление, обеспечивающее её достижение. В первой системе
$s=x_1+x_2+\sin x_1$, во второй $s=3x_1+x_2$; эти поверхности выбраны
из разных первых уравнений объектов.

\item \textbf{Что означает $\beta_0>0$?}
Это строгий запас над оценкой неопределённости:
$\beta\ge\varrho+\beta_0$. Он даёт
$\frac{d}{dt}|s|\le-g_0\beta_0$ и конечное время достижения.
Если убрать запас, простого строгого неравенства для конечного времени
уже может не быть.

\item \textbf{Почему нельзя реализовать идеальное скольжение простой подстановкой $\sgn0=0$?}
В нуле поверхности требуется эквивалентный вход, компенсирующий дрейф.
В модели Филиппова он выбирается из интервала значений реле.
Обычное нулевое значение знака в одной точке этого не описывает.

\item \textbf{Почему непрерывность не означает дифференцируемость?}
Функция $\sat(s/\varepsilon)$ непрерывна, но её производная меняется скачком
при $|s|=\varepsilon$. Регулятор непрерывный; называть его обязательно
гладким в смысле $C^1$ было бы неверно.

\item \textbf{Почему насыщенный регулятор системы 1 приводит точно к нулю?}
После достижения полосы отдельная полная функция
$V_c=(x_1^2+s^2)/2$ удовлетворяет $\dot V_c\le-V_c$ при
$0<\varepsilon\le1/4$. Неопределённость исчезает в нуле,
а выбранное усиление перекрывает смешанные члены.
Это специальный результат данной системы.

\item \textbf{Почему для маятника остаётся окрестность?}
При ненулевом неизвестном $h(t)$ в самом нуле требуется неизвестный
момент $-mlh(t)$. Непрерывный статический закон насыщения в нуле даёт
единственное значение $0$. Поэтому доказаны оценки
$\limsup|\theta|\le\varepsilon$ и
$\limsup|\dot\theta|\le2\varepsilon$.

\item \textbf{Почему усиление адаптации не обязательно сходится к неизвестному параметру?}
Для $\dot c=\gamma x^2$ при $x=0$ изменение $c$ прекращается при любом его
значении. Функция Ляпунова доказывает $x\to0$ и ограниченность $c$,
а не идентификацию $\Delta a$.
\end{enumerate}

\subsection{Как объяснить численное моделирование}
\addition В работах 4 и 5 используется метод Рунге-Кутты четвёртого порядка.
Для автономной правой части $F$ и шага $h$ один шаг имеет вид
\[
 \begin{aligned}
 k_1&=F(x_i),& k_2&=F(x_i+hk_1/2),\\
 k_3&=F(x_i+hk_2/2),& k_4&=F(x_i+hk_3),\\
 x_{i+1}&=x_i+\frac h6(k_1+2k_2+2k_3+k_4).
 \end{aligned}
\]
В работе 4 заданы $h=0{,}001$ с, $t\in[0,10]$ с,
$x(0)=(1;-0{,}5)^T$ для первых двух систем и
$x(0)=(0{,}5;1;-0{,}5;0{,}25)^T$ для третьей.
Правая часть на каждом этапе вычисляется вместе с обратной связью
в промежуточном состоянии. Для третьей системы отображаются
$x_1$, $x_2-1$, $x_3$, $x_4$, а не требование занулить все исходные координаты.

В работе 5 используются тот же шаг, интервал $[0,12]$ с,
параметры $\theta_1=1$, $\theta_2=-1$, $a_1=a_2=2$ и начальные
состояния $(1,0)^T$, $(1{,}5;0)^T$ соответственно. Эти числа известны
модели объекта, а робастный закон использует только границы и состояние.
Для непрерывного закона взято $\varepsilon=0{,}1$.

Разрывная динамика интегрируется до пересечения $s=0$.
В сокращённом скрипте момент пересечения внутри шага оценивается линейной
интерполяцией значения $s$:
\[
 \tau\approx h\frac{s_i}{s_i-s_{i+1}},\qquad 0\le\tau\le h.
\]
После него интегрируется движение на поверхности и отображается
эквивалентное управление. Для непрерывного закона насыщение вычисляется
непосредственно на каждом этапе. Поэтому отсутствие частых переключений
на идеальном графике реле не является доказательством отсутствия chattering
у реального привода.

Один численный опыт подтверждает поведение выбранной модели при конкретных
начальных условиях и параметрах. Гарантия для всех допустимых параметров
и всех начальных состояний обеспечивается аналитическими оценками,
а не самим видом графика.
